\section{Introduction}
Language-model research systems increasingly retrieve sources, synthesize claims, attach citations, run automated verification, and decide whether outputs may be acted upon. This creates two coupled scientific problems.  First, a claim may be correct but attributed to the wrong source; a source may support a nearby claim but not the assigned one; a checker may be non-empty yet non-discriminating or share lineage with the answer; and an evaluator may be contaminated or tampered with after the fact. In each case, a scalar confidence score can remain high even though a prerequisite for scientific authority is absent.  Second, a benchmark may report a clean number without measuring the intended competence: its gold terminal may be exposed by construction artifacts, unavailable in a comparator's interface, or constant over the entire panel.

We study a stricter object: a \emph{scientific-authority transition}. A proposal can be plausible and well supported but is promoted only when a registered set of content, provenance, checker, evaluator, scope, dependence, and freshness obligations all pass. Unresolved prerequisites do not contribute a low score; they block escalation. This design is motivated by source-aware verification, multi-source attribution, claim-level auditability, evidence influence, iterative fact checking, evaluator-tampering benchmarks, search-time contamination, partial-evidence awareness, and evidence-backed authorization \citep{provenanceguard2026,attributionbench2024,aar2026,provenai2026,fire2024,rewardhacking2026,stc2026,deepsciverify2026,partial2026,authgraph2026,fava2026,cva2026}. Recent academic-integrity and agentic-abstention benchmarks also make explicit that declining to complete an unsupported action can itself be the correct behavior rather than a failure to optimize task success \citep{sciintegrity2026,agentabstain2026}.

The contribution moves above those component techniques rather than relabeling them.  We ask when a verification axis can support a scientific claim at all, and when donor-complete local mechanisms jointly suffice for target scientific promotion.  Section~\ref{sec:axis-theory} gives an exact attainability criterion, a mixed-fibre lower bound, a donor-product factorization theorem, and a diagnostic separation among construction saturation, nuisance leakage, interface limitation, and decision failure.  The empirical studies then ask whether composing the donor mechanisms as a \emph{non-compensatory, protected authority transition} reduces false scientific promotion while preserving clean promotion. Throughout, the pipeline built to that rule is called the \emph{governed pipeline}, so that each claim reads as a claim about the rule rather than about a product; Section~\ref{sec:methods} names the implementation that instantiates it.

The V2 study was executed after the subject, comparator/ablation mechanisms, evaluator, statistical rules, and split-generation code were frozen. A first 420-case protected campaign on an older subject remains valid evidence for that older revision but does not authorize the repaired one. A separate 39-case live-model arm remains exploratory because post-run audit found label/denominator inconsistencies and a hidden-family execution leak. We therefore generated a fresh post-repair hidden split and executed the publication-authorizing V2 campaign without reusing either prior outcome for tuning.  V3 and P4-X are separately identified successors; neither is pooled with V2 or used to rewrite its terminals.

Our contributions are:
\begin{enumerate}
  \item a general finite-space theory of verification-axis attainability, including the exact Bayes-risk expression, a necessary-and-sufficient zero-risk condition, a total-variation bound for wider claim identifiability, a nuisance-identifiability criterion, and a donor-product factorization boundary;
  \item a protected V2 result of 0/360 versus 180/360 false promotions with 60/60 clean promotion in both arms, together with the retained negative finding that H3 was non-discriminating because its construction was saturated;
  \item an exact-axis V3 repair that clears the frozen shortcut probes and obtains 30/30 correct \textsc{CannotCheck} terminals versus 0/30 and 15/30 for the relevant comparators, while explicitly identifying that result as interface expressiveness rather than general epistemic judgement; and
  \item P4-X, a donor-engulfing successor in which the non-compensatory scientific-promotion relation scores 400/400 versus 250/400 for a strong generic-authorization product and 50/400 for a compensatory product, while an information-equivalent typed product ties at 400/400 as the factorization theorem requires.
\end{enumerate}

The widest empirical target---naturalistic, cross-domain scientific-promotion judgement against \textsc{CannotCheck}-capable external comparators---remains prospective.  Its frozen design contains 768 paired source-family clusters across four domains and eight unidentifiability mechanisms, with an independently held external panel and source-disjoint replication required before any such claim can be promoted.
